\documentclass[11pt]{article}

% ============================================================
% GA-TL-01 — Identificación y clasificación de tipos de requisitos
% Ingeniería de Requerimientos (ISR-401) · UTEQ · SPA 2026-2027
% Compilación: pdflatex -> bibtex -> pdflatex -> pdflatex
% ============================================================

\usepackage[utf8]{inputenc}
\usepackage[T1]{fontenc}
% babel no disponible en este entorno
\usepackage[scaled]{helvet}
\renewcommand{\familydefault}{\sfdefault}
\usepackage[a4paper,top=2.6cm,bottom=2.3cm,left=2cm,right=2cm,headheight=34pt]{geometry}
\usepackage{graphicx}
\usepackage{array}
\usepackage{tabularx}
\usepackage{longtable}
\usepackage{pdflscape}
\usepackage{colortbl}
\usepackage{xcolor}
\usepackage{enumitem}
\usepackage{fancyhdr}
\usepackage{lastpage}
\usepackage[most]{tcolorbox}
\usepackage{amsmath}
\usepackage{amssymb}
\usepackage[numbers,sort&compress]{natbib}
\usepackage{url}
\usepackage{needspace}
\usepackage{caption}
\usepackage[hidelinks]{hyperref}

% ---------- Textos en español (sin babel) ----------
\renewcommand{\tablename}{Tabla}
\renewcommand{\refname}{Referencias}

% ---------- Paleta institucional ----------
\definecolor{utegreen}{HTML}{0B7A3B}
\definecolor{utegreendk}{HTML}{075C2C}
\definecolor{midblue}{HTML}{1F5C8B}
\definecolor{gold}{HTML}{D4A017}
\definecolor{rowgray}{HTML}{EFEFEF}
\definecolor{boxgray}{HTML}{F4F6F4}
\definecolor{lightred}{HTML}{FBEAEA}
\definecolor{lightgold}{HTML}{FBF3DC}
\definecolor{lightblue}{HTML}{EAF1F7}
\definecolor{lightgreen}{HTML}{DCEFE3}

% ---------- Columnas personalizadas ----------
\newcolumntype{L}[1]{>{\raggedright\arraybackslash}p{#1}}
\newcolumntype{C}[1]{>{\centering\arraybackslash}p{#1}}
\newcolumntype{Y}{>{\raggedright\arraybackslash}X}

% ---------- Encabezado y pie ----------
\pagestyle{fancy}
\fancyhf{}
\lhead{\footnotesize UTEQ \,$\cdot$\, Ingeniería de Requerimientos (ISR-401)}
\rhead{\footnotesize GA-TL-01 \,$\cdot$\, SPA 2026--2027}
\lfoot{\footnotesize Prof. Gleiston C. Guerrero-Ulloa, PhD}
\rfoot{\footnotesize Página \thepage\ de \pageref{LastPage}}
\renewcommand{\headrulewidth}{0.6pt}
\renewcommand{\footrulewidth}{0.4pt}

\setlength{\parindent}{0pt}
\setlength{\parskip}{5pt}
\setlength{\emergencystretch}{3em}
\renewcommand{\arraystretch}{1.25}

% ---------- Atajos ----------
\newcommand{\thdr}[1]{\cellcolor{utegreen}\textcolor{white}{\textbf{#1}}}
\newcommand{\seccion}[1]{\section{\textcolor{utegreendk}{#1}}}
\newcommand{\subseccion}[1]{\subsection{\textcolor{midblue}{#1}}}
\renewcommand{\thesection}{\arabic{section}.}
\renewcommand{\thesubsection}{\arabic{section}.\arabic{subsection}.}

\newtcolorbox{cajaclave}[1][]{colback=boxgray,colframe=utegreen,boxrule=0.8pt,arc=2mm,fonttitle=\bfseries,breakable,#1}
\newtcolorbox{cajapiso}[1][]{colback=lightred,colframe=red!60!black,boxrule=1pt,arc=2mm,fonttitle=\bfseries,breakable,#1}
\newtcolorbox{cajaaviso}[1][]{colback=lightgold,colframe=gold!80!black,boxrule=0.8pt,arc=2mm,fonttitle=\bfseries,breakable,#1}

\begin{document}

\noindent\includegraphics[width=\linewidth]{logoband.png}

\vspace{4pt}
\begin{center}
	{\small Facultad de Ciencias de la Computación \,\textcolor{gold}{$\bullet$}\, Carrera de Software \,\textcolor{gold}{$\bullet$}\, Ingeniería de Requerimientos (ISR-401)}\\[8pt]
	{\LARGE\textbf{\textcolor{utegreendk}{GUÍA Y RÚBRICA DE EVALUACIÓN}}}\\[4pt]
	{\large\textbf{\textcolor{midblue}{ACTIVIDAD FORMATIVA}}}\\[6pt]
	{\Large\textbf{GA-TL-01}}\\[2pt]
	{\large Identificación y clasificación de tipos de requisitos}
\end{center}

\vspace{2pt}
\noindent\textcolor{gold}{\rule{\linewidth}{1.5pt}}

\seccion{Datos generales}

\noindent\begin{tabularx}{\linewidth}{|L{0.30\linewidth}|Y|}
	\hline
	\thdr{Campo} & \thdr{Detalle}\\
	\hline
	Asignatura & Ingeniería de Requerimientos (ISR-401) \,$\cdot$\, 4.\textsuperscript{o} nivel \,$\cdot$\, Paralelo A\\
	\hline
	\rowcolor{rowgray} Docente & Gleiston Cicerón Guerrero Ulloa, PhD \,$\cdot$\, \texttt{gguerrero@uteq.edu.ec}\\
	\hline
	Unidad & Unidad 1. Fundamentos de Ingeniería de Requisitos y Proceso de Requisitos\\
	\hline
	\rowcolor{rowgray} Resultado de aprendizaje & Aplica los fundamentos teóricos de la IR identificando tipos, propiedades y el proceso de requisitos para analizar su aplicabilidad en proyectos de software reales.\\
	\hline
	Carácter & Formativo: sin calificación; su entrega válida es requisito de la sumativa vinculada\\
	\hline
	\rowcolor{rowgray} Prepara (sumativas vinculadas) & GA-PE-U, EV-PRAC-01\\
	\hline
	Modalidad & Individual\\
	\hline
	\rowcolor{rowgray} Semana & 1 (26 de octubre -- 1 de noviembre de 2026)\\
	\hline
	Publicación de la guía & lunes 26 de octubre de 2026\\
	\hline
	\rowcolor{rowgray} Cierre & sábado 31 de octubre de 2026, a las 23:59\\
	\hline
	Repositorio & Repositorio personal \texttt{ISR401-<ApellidoNombre>}\\
	\hline
\end{tabularx}

\seccion{Propósito}

Practicar en clase la distinción entre tipos de requisitos y la formulación de requisitos de calidad cuantificados, antes de hacerlo en equipo en la práctica experimental de la unidad 1 (GA-PE-U) y de forma individual en la evaluación EV-PRAC-01.

\seccion{Fundamento}

La distinción entre requisitos funcionales, de calidad y restricciones, entre requisitos de producto y de proceso y entre requisitos del sistema y del software es la base de toda la IR \citep{pohl_2025,iso29148_2018}. Los requisitos de calidad se clasifican con ISO/\allowbreak{}IEC 25010:2023 \citep{iso25010_2023}.

\seccion{Desarrollo de la actividad}

\begin{enumerate}[leftmargin=20pt,itemsep=4pt]
	\item \textbf{Enunciados.} Lea el caso de la Plataforma de gestión de citas médicas (Anexo A) y use los 20 enunciados de la Tabla~\ref{tab:enunciados}.
	\item \textbf{Clasificación.} Clasifique cada enunciado (funcional, calidad con característica ISO/\allowbreak{}IEC 25010, restricción; producto o proceso; sistema o software) y justifique en una línea.
	\item \textbf{Cuantificación.} Reescriba 5 requisitos de calidad vagos («rápido», «seguro», «fácil de usar») como requisitos cuantificados con métrica y umbral.
	\item \textbf{Discusión.} Compare sus respuestas con un compañero y anote las discrepancias y su resolución.
\end{enumerate}

\seccion{Productos entregables}

{\small
\begin{longtable}{|L{0.27\linewidth}|L{0.43\linewidth}|L{0.21\linewidth}|}
	\caption{Productos que deben estar en el repositorio al cierre.}\\
	\hline
	\thdr{Ruta en el repositorio} & \thdr{Contenido} & \thdr{Mínimo exigido}\\
	\hline
	\endfirsthead
	\hline
	\thdr{Ruta en el repositorio} & \thdr{Contenido} & \thdr{Mínimo exigido}\\
	\hline
	\endhead
	\texttt{GA-TL-01/\allowbreak{}main.tex y main.pdf} & Tabla de clasificación, 5 requisitos cuantificados y discrepancias & 20 + 5\\
	\hline
\end{longtable}
}

\seccion{Normas de entrega}

\begin{enumerate}[leftmargin=20pt,itemsep=3pt]
	\item \textbf{Repositorio.} Todo el trabajo se desarrolla en el repositorio personal \texttt{ISR401-<ApellidoNombre>}, que debe ser público en GitHub. Allí deben estar los fuentes (\texttt{.tex}, \texttt{.bib}, figuras, fuentes editables de los modelos y evidencias) y el PDF compilado.
	\item \textbf{Documento en el SGA.} En la tarea del SGA se sube solamente un PDF de \textbf{una sola página} con los datos informativos: estudiante o estudiantes, código y nombre de la tarea, fecha límite de entrega y la URL del repositorio escrita completa en una sola línea. No se sube ningún otro archivo.
	\item \textbf{\LaTeX{} reproducible.} El PDF debe poder regenerarse clonando el repositorio y compilando el archivo \texttt{.tex}. El \texttt{README.md} documenta compilador, archivo principal, orden de comandos y dependencias.
	\item \textbf{Referencias.} Estilo IEEE. Toda referencia debe existir y sus datos (autores, año, título, revista o editorial, volumen, número, páginas y DOI o ISBN) deben corresponder exactamente a la obra real.
	\item \textbf{Plazo.} El cierre es el sábado 31 de octubre de 2026, a las 23:59. Se evalúa el contenido del repositorio hasta ese momento; lo agregado después no se considera.
	\item \textbf{Trabajo anticipado.} El trabajo se realiza durante la semana, no a última hora; el historial de commits debe mostrarlo. La lista de cotejo registra la constancia observada.
	\item \textbf{Datos personales.} Ningún dato personal de terceros (cédulas, teléfonos, correos personales, audios o videos de participantes, consentimientos firmados) se publica en el repositorio. Esas evidencias se guardan en una carpeta institucional privada (cuenta \texttt{@uteq.edu.ec}) compartida con el docente, y en el repositorio se registra un manifiesto con nombre de archivo, fecha, duración y huella SHA-256 \citep{lopdp_2021}.
	\item \textbf{Uso de IA como apoyo.} Si se usa una herramienta de IA solo como apoyo (ortografía, error de compilación), se declara en una sección del documento: herramienta, versión, propósito y cómo se verificó lo obtenido.
\end{enumerate}

\seccion{Condiciones de entrega válida}

\begin{cajapiso}[title={Condiciones de entrega válida: su incumplimiento equivale a no entregar}]
	\begin{enumerate}[label=\textbf{P\arabic*.},leftmargin=24pt,itemsep=3pt]
		\item \textbf{PDF de una sola página con datos informativos y URL.} En el SGA se sube solamente un PDF de UNA SOLA página con los datos informativos (estudiante o estudiantes, tarea, fecha límite de entrega) y, en una sola línea, la URL del repositorio. Si el PDF tiene más de una página, si se suben otros archivos, si la URL no cabe en una sola línea o está rota, o si el repositorio no es accesible, la actividad formativa se considera \textbf{NO ENTREGADA} y el estudiante no puede presentar la sumativa vinculada, en la que obtiene CERO.
		\item \textbf{PDF reproducible.} Si el PDF no se genera correctamente a partir del \LaTeX{} al clonar el repositorio y compilar el archivo \texttt{.tex} siguiendo el \texttt{README.md} (compilador, orden de comandos, archivo principal y dependencias), la actividad formativa se considera \textbf{NO ENTREGADA} y el estudiante no puede presentar la sumativa vinculada, en la que obtiene CERO.
		\item \textbf{Plazo.} Si el contenido evaluable se sube después de la fecha y hora de cierre, la actividad formativa se considera \textbf{NO ENTREGADA} y el estudiante no puede presentar la sumativa vinculada, en la que obtiene CERO.
		\item \textbf{Integridad.} Plagio verificado, cualquiera de los indicios graves G1 a G6 de uso de IA, o contenido generado por IA presentado como propio: la actividad formativa se considera \textbf{NO ENTREGADA} y el estudiante no puede presentar la sumativa vinculada, en la que obtiene CERO, sin perjuicio de las acciones disciplinarias que correspondan.
	\end{enumerate}
\end{cajapiso}

\seccion{Integridad académica e indicios de uso de IA}

Todo trabajo debe ser de autoría propia. Se aplica el Código de Ética de la ACM \citep{acm_code_2018} y la sección 6.1 del encuadre.

\begin{cajapiso}[title={Indicios graves de uso de IA: cada uno, por sí solo, implica CERO (P4)}]
	\begin{enumerate}[label=\textbf{G\arabic*.},leftmargin=24pt,itemsep=3pt]
		\item \textbf{Referencias con datos que no corresponden a la obra real:} DOI que lleva a una obra con autores, título, año, revista, volumen o páginas diferentes; DOI, ISBN o URL inexistentes; obras que no existen.
		\item \textbf{Citas que la fuente no respalda:} afirmaciones, cifras o citas textuales atribuidas a una fuente que no las contiene, o cláusulas, secciones o tablas de una norma (IEEE 29148, ISO/\allowbreak{}IEC 25010) que no existen o no dicen lo afirmado.
		\item \textbf{Restos de un asistente de IA} en el documento, el código, los comentarios o los commits (por ejemplo, «Claro, aquí tienes\ldots», «Como modelo de lenguaje\ldots»).
		\item \textbf{Evidencias de campo fabricadas:} entrevistas, encuestas, participantes, transcripciones, actas o fotografías que no existen o no tienen respaldo original.
		\item \textbf{Resultados o artefactos inexistentes:} el texto describe modelos, matrices, pruebas o métricas que no están en el repositorio o no coinciden con él.
		\item \textbf{Incapacidad de explicar lo entregado:} el estudiante no puede justificar ni modificar, cuando se le pide, lo que escribió, modeló o especificó.
	\end{enumerate}
\end{cajapiso}

\begin{cajaaviso}[title={Indicios complementarios (se acumulan)}]
	\begin{enumerate}[label=\textbf{I\arabic*.},leftmargin=24pt,itemsep=3pt]
		\item Redacción uniforme de principio a fin, o estilo distinto al del estudiante en otras entregas o en su exposición oral.
		\item Uso excesivo de los mismos conectores y fórmulas («además», «cabe destacar», «es importante señalar», «en conclusión»).
		\item Contenido genérico: requisitos, actores o reglas aplicables a cualquier sistema, sin relación con lo elicitado.
		\item Entrega en uno o pocos commits masivos, o mensajes de commit idénticos o automáticos.
		\item Residuos de formato ajeno a \LaTeX{} (marcas de Markdown, emojis) o metadatos del PDF de otra herramienta.
		\item Estructura, redacción o errores muy similares entre estudiantes o equipos distintos.
		\item Traducción literal o mezcla de idiomas sin justificación.
	\end{enumerate}
	Con dos indicios complementarios se resta 1,5 puntos (nivel B) y con tres o más se restan 3,0 (nivel C). La declaración de uso de IA no exime de responsabilidad ni anula estos indicios.
\end{cajaaviso}

\seccion{Lista de cotejo para la retroalimentación}

La actividad no tiene calificación. El docente marca cada indicador y escribe una observación que el estudiante debe atender en la actividad sumativa vinculada. \textbf{L} = Logrado (cumple con la calidad descrita); \textbf{EP} = En proceso (presente pero incompleto, impreciso o con errores que se señalan); \textbf{NL} = No logrado (ausente, incorrecto o no verificable). Que un elemento esté presente no basta: se valora su calidad.

{\small
\begin{longtable}{|C{0.04\linewidth}|L{0.47\linewidth}|C{0.035\linewidth}|C{0.035\linewidth}|C{0.035\linewidth}|L{0.17\linewidth}|}
	\caption{Lista de cotejo de GA-TL-01.}\\
	\hline
	\thdr{N.\textsuperscript{o}} & \thdr{Indicador y calidad esperada (Logrado)} & \thdr{L} & \thdr{EP} & \thdr{NL} & \thdr{Observación}\\
	\hline
	\endfirsthead
	\hline
	\thdr{N.\textsuperscript{o}} & \thdr{Indicador y calidad esperada (Logrado)} & \thdr{L} & \thdr{EP} & \thdr{NL} & \thdr{Observación}\\
	\hline
	\endhead
	1 & Los 20 enunciados clasificados en las tres dimensiones con justificación. & $\square$ & $\square$ & $\square$ & \\
	\hline
	2 & Características y subcaracterísticas de ISO/\allowbreak{}IEC 25010:2023 asignadas correctamente. & $\square$ & $\square$ & $\square$ & \\
	\hline
	3 & Los 5 requisitos de calidad reescritos con métrica, unidad y umbral verificables. & $\square$ & $\square$ & $\square$ & \\
	\hline
	4 & Discrepancias con el compañero registradas y resueltas con argumento. & $\square$ & $\square$ & $\square$ & \\
	\hline
	5 & Documento propio, compilado y en el repositorio personal. & $\square$ & $\square$ & $\square$ & \\
	\hline
\end{longtable}
}

\begin{cajaclave}[title={Cómo se usa esta retroalimentación}]
	El docente devuelve la lista de cotejo en el SGA antes de la sesión siguiente. Cada indicador marcado como EP o NL debe corregirse antes de presentar GA-PE-U y EV-PRAC-01; el docente verificará esa corrección al calificarla. Entregar esta formativa de forma válida es requisito para presentar GA-PE-U y EV-PRAC-01.
\end{cajaclave}

% ============================================================
% ANEXOS: material que entrega el docente
% ============================================================
\clearpage
\seccion{Anexo A. Caso: Plataforma de gestión de citas médicas}
\label{sec:caso}

\subseccion{Contexto de la organización}
El Policlínico Ribera del Quevedo (organización ficticia creada para esta actividad) es un centro privado de consulta externa ubicado en Quevedo, provincia de Los Ríos. Atiende de lunes a sábado, de 07:00 a 19:00, con 24 médicos distribuidos en 9 especialidades (medicina general, pediatría, ginecología, cardiología, traumatología, dermatología, endocrinología, gastroenterología y oftalmología) y 18 consultorios. Tiene alrededor de 12\,000 pacientes registrados y agenda en promedio 450 citas por día. Cerca del 35\,\% de sus pacientes tiene más de 60 años y muchos de ellos acuden acompañados de un familiar que les ayuda con los trámites.

\subseccion{Situación actual y problema}
Hoy las citas se agendan por teléfono o de forma presencial en recepción. Tres recepcionistas registran las citas en una hoja de cálculo compartida y, cuando la red falla, en un cuaderno que luego transcriben. Entre las 07:00 y las 09:00 las líneas telefónicas se saturan: muchos pacientes reportan que llaman varias veces sin ser atendidos y terminan acudiendo en persona para hacer fila desde la madrugada. La dirección médica estima, con sus propios registros, que el 22\,\% de los pacientes no se presenta a su cita sin avisar, de modo que quedan horas de consulta vacías mientras otros pacientes esperan días por un turno. Además, se han detectado citas duplicadas (dos pacientes asignados al mismo médico a la misma hora) por errores al copiar datos entre la hoja de cálculo y el cuaderno.

El policlínico ya cuenta con un sistema de escritorio de facturación electrónica, adquirido a un proveedor externo, que no se va a reemplazar; el proveedor ofrece una interfaz de servicios web documentada para registrar los cobros. El laboratorio clínico usa su propio sistema y no forma parte de este proyecto. Las historias clínicas siguen en papel y su digitalización se hará en un proyecto posterior.

\subseccion{Interesados}
\begin{itemize}
	\item \textbf{Pacientes y sus familiares:} quieren reservar, reprogramar o cancelar citas sin llamar ni hacer fila, y recibir recordatorios.
	\item \textbf{Recepcionistas:} quieren dejar de transcribir datos, registrar la llegada del paciente y atender a quienes no usan internet.
	\item \textbf{Médicos:} quieren conocer su agenda del día, saber qué pacientes llegaron y marcar las consultas atendidas.
	\item \textbf{Director médico:} quiere reducir el ausentismo, aprovechar mejor las horas de consulta y recibir reportes por especialidad.
	\item \textbf{Administrador del policlínico:} quiere gestionar médicos, horarios, especialidades y consultorios, y controlar los cobros.
	\item \textbf{Área de TI (dos personas):} mantendrá la plataforma una vez entregada; domina Java y PostgreSQL.
	\item \textbf{Proveedor del sistema de facturación:} define cómo se integra cualquier sistema nuevo con el suyo.
\end{itemize}

\subseccion{Lo que se espera de la plataforma}
El director médico solicitó una plataforma web, acompañada de una aplicación móvil, que permita al paciente consultar la disponibilidad por especialidad y médico, reservar una cita, reprogramarla o cancelarla, y recibir recordatorios por mensaje de texto y correo electrónico. Cuando no haya turnos libres, el paciente podrá inscribirse en una lista de espera y ser avisado si se libera un horario. En la sala de espera se instalará un quiosco con pantalla táctil e impresora térmica donde el paciente confirmará su llegada y recibirá un ticket con su número de turno. El médico verá en su consultorio la lista de pacientes del día y su estado. El administrador configurará horarios, ausencias de los médicos y feriados, y el director médico obtendrá reportes mensuales de citas atendidas, canceladas e inasistencias.

Las reglas que hoy aplica la recepción y que deben mantenerse son: una cita dura 20 minutos, salvo en cardiología y traumatología, donde dura 30; un paciente no puede tener dos citas pendientes con la misma especialidad; una cita cancelada con menos de cuatro horas de anticipación se considera inasistencia; y un paciente con tres inasistencias en un mismo mes solo puede agendar presencialmente durante el mes siguiente.

\subseccion{Condiciones del proyecto}
El proyecto tiene un plazo de seis meses y se desarrollará con un enfoque iterativo, con entregas cada dos semanas revisadas por el director médico. Los datos que manejará la plataforma incluyen información de salud de los pacientes, que la Ley Orgánica de Protección de Datos Personales del Ecuador considera datos sensibles. El director médico pidió expresamente que la plataforma «sea rápida, segura, fácil de usar para los adultos mayores, esté siempre disponible y se integre sin problemas con la facturación», pero no supo precisar qué significa cada una de esas expresiones. Como referencia para cuantificarlas, el área de TI informa que los lunes a las 07:00 se espera el mayor número de usuarios conectados al mismo tiempo, alrededor de 300, y que la mayoría de los pacientes usa teléfonos Android de gama media o baja.

\subseccion{Fuera del alcance}
Quedan fuera de este proyecto la historia clínica electrónica, la telemedicina, la farmacia, el sistema del laboratorio y la emisión de facturas, que seguirá a cargo del sistema actual.

\subseccion{Enunciados para clasificar}
La Tabla~\ref{tab:enunciados} contiene los 20 enunciados que usted debe clasificar en las tres dimensiones indicadas en la sección «Desarrollo de la actividad». Entre ellos hay cinco requisitos de calidad redactados de forma vaga: identifíquelos y reescríbalos con métrica, unidad y umbral, apoyándose en los datos de este caso.

{\small
\begin{longtable}{|>{\centering\arraybackslash}p{1.5cm}|p{\dimexpr\textwidth-1.5cm-4\tabcolsep-3\arrayrulewidth\relax}|}
	\caption{Enunciados de la Plataforma de gestión de citas médicas.}\label{tab:enunciados}\\
	\hline
	\thdr{Código} & \thdr{Enunciado} \\
	\hline
	\endfirsthead
	\hline
	\thdr{Código} & \thdr{Enunciado} \\
	\hline
	\endhead
	E01 & El paciente podrá reservar una cita eligiendo especialidad, médico, fecha y hora entre los horarios disponibles. \\ \hline
	E02 & La plataforma enviará al paciente un recordatorio por mensaje de texto 24 horas antes de su cita. \\ \hline
	E03 & La aplicación debe ser rápida. \\ \hline
	E04 & El recepcionista podrá registrar la llegada del paciente que se acerca a la ventanilla. \\ \hline
	E05 & La plataforma se desarrollará en Java 21 con base de datos PostgreSQL, tecnologías que domina el área de TI. \\ \hline
	E06 & El médico podrá consultar la lista de pacientes del día y marcar cada consulta como atendida. \\ \hline
	E07 & Los datos de salud de los pacientes deben estar seguros. \\ \hline
	E08 & La aplicación móvil debe funcionar en teléfonos con Android 10 o superior. \\ \hline
	E09 & Una cita cancelada con menos de cuatro horas de anticipación se registrará como inasistencia. \\ \hline
	E10 & La plataforma debe ser fácil de usar para los pacientes adultos mayores. \\ \hline
	E11 & El proyecto debe entregarse en un plazo máximo de seis meses. \\ \hline
	E12 & La sesión de un usuario se cerrará automáticamente tras 10 minutos sin actividad. \\ \hline
	E13 & El quiosco de la sala de espera leerá el número de cédula ingresado en su pantalla táctil e imprimirá un ticket con el número de turno. \\ \hline
	E14 & La plataforma debe estar siempre disponible. \\ \hline
	E15 & El director médico podrá generar un reporte mensual de inasistencias por especialidad. \\ \hline
	E16 & Cada incremento se validará con el director médico al final de cada iteración de dos semanas. \\ \hline
	E17 & El tratamiento de los datos personales de los pacientes debe cumplir la Ley Orgánica de Protección de Datos Personales del Ecuador. \\ \hline
	E18 & La plataforma debe integrarse sin problemas con el sistema de facturación actual. \\ \hline
	E19 & Si el servidor principal falla, la plataforma debe restablecer el servicio en menos de 30 minutos sin perder ninguna cita confirmada. \\ \hline
	E20 & Todo cambio en el código fuente será revisado por otro integrante del equipo antes de integrarse a la rama principal. \\ \hline
\end{longtable}
}

\bibliographystyle{unsrtnat}
\bibliography{referencias}

\end{document}
